% GENERATED FILE. Edit canonical lessons, then run npm run generate:books.
% canonical-source-hash: fnv1a64:34d368581710803b
% canonical-lessons: TE-C175-gumpu, TE-C175-prapancam, TE-C175-bhukampam, TE-C175-dummu, TE-C175-poga

\chapter{Crowd, World, Earthquake, Dust, Smoke}
\label{ch:te-crowd-world-earthquake-dust-smoke}
\hlchaptermodality{175}

\begin{chapteropening}
\textbf{By the end of this chapter:} \emph{I can say a crowd, the world, an earthquake, dust and smoke.}

Complete the last lesson of chapter 175: I can say a crowd, the world, an earthquake, dust and smoke.
\end{chapteropening}

\section[\texorpdfstring{\te{గుంపు} (gumpu)}{గుంపు (gumpu)}]{\te{గుంపు} (gumpu) — a crowd}

\label{lesson:TE-C175-gumpu}



\begin{quote}
Before the new one: say the Telugu for a cave, then the Telugu for a valley.
\end{quote}

\subsection*{You'll want to know: \te{గుంపు}}
\textbf{\te{గుంపు}} — \emph{gumpu} — \textquotedblleft{}a crowd\textquotedblright{}.

\begin{grammarlens}[title={What you've built}]
One more word for this chapter.
\end{grammarlens}

\subsection*{Your turn}
\begin{itemize}
\raggedright
  \item \emph{Say it:} \emph{gumpu}
  \item \emph{Say it:} \emph{gumpu}, once more
  \item \emph{Say it:} say \emph{gumpu}
  \item \emph{Recall:} say the Telugu for a cave, then the Telugu for a valley, then say \emph{gumpu} again
\end{itemize}

\begin{tcolorbox}[breakable,colback=teal!4,colframe=teal!35!black,title={Before you move on}]
What does \te{గుంపు} mean? (\textquotedblleft{}A crowd\textquotedblright{}.) Say it once more.
\end{tcolorbox}
\section[\texorpdfstring{\te{ప్రపంచం} (prapañcaṁ)}{ప్రపంచం (prapañcaṁ)}]{\te{ప్రపంచం} (prapañcaṁ) — the world}

\label{lesson:TE-C175-prapancam}



\begin{quote}
Before the new one: say the Telugu for a valley, then the Telugu for a crowd.
\end{quote}

\subsection*{You'll want to know: \te{ప్రపంచం}}
\textbf{\te{ప్రపంచం}} — \emph{prapañcaṁ} — \textquotedblleft{}the world\textquotedblright{}.

\begin{grammarlens}[title={What you've built}]
One more word for this chapter.
\end{grammarlens}

\subsection*{Your turn}
\begin{itemize}
\raggedright
  \item \emph{Say it:} \emph{prapañcaṁ}
  \item \emph{Say it:} \emph{prapañcaṁ}, once more
  \item \emph{Say it:} say \emph{prapañcaṁ}
  \item \emph{Recall:} say the Telugu for a valley, then the Telugu for a crowd, then say \emph{prapañcaṁ} again
\end{itemize}

\begin{tcolorbox}[breakable,colback=teal!4,colframe=teal!35!black,title={Before you move on}]
What does \te{ప్రపంచం} mean? (\textquotedblleft{}The world\textquotedblright{}.) Say it once more.
\end{tcolorbox}
\section[\texorpdfstring{\te{భూకంపం} (bhūkampaṁ)}{భూకంపం (bhūkampaṁ)}]{\te{భూకంపం} (bhūkampaṁ) — an earthquake}

\label{lesson:TE-C175-bhukampam}



\begin{quote}
Before the new one: say the Telugu for a crowd, then the Telugu for the world.
\end{quote}

\subsection*{You'll want to know: \te{భూకంపం}}
\textbf{\te{భూకంపం}} — \emph{bhūkampaṁ} — \textquotedblleft{}an earthquake\textquotedblright{}.

\begin{grammarlens}[title={What you've built}]
One more word for this chapter.
\end{grammarlens}

\subsection*{Your turn}
\begin{itemize}
\raggedright
  \item \emph{Say it:} \emph{bhūkampaṁ}
  \item \emph{Say it:} \emph{bhūkampaṁ}, once more
  \item \emph{Say it:} say \emph{bhūkampaṁ}
  \item \emph{Recall:} say the Telugu for a crowd, then the Telugu for the world, then say \emph{bhūkampaṁ} again
\end{itemize}

\begin{tcolorbox}[breakable,colback=teal!4,colframe=teal!35!black,title={Before you move on}]
What does \te{భూకంపం} mean? (\textquotedblleft{}An earthquake\textquotedblright{}.) Say it once more.
\end{tcolorbox}
\section[\texorpdfstring{\te{దుమ్ము} (dummu)}{దుమ్ము (dummu)}]{\te{దుమ్ము} (dummu) — dust}

\label{lesson:TE-C175-dummu}



\begin{quote}
Before the new one: say the Telugu for the world, then the Telugu for an earthquake.
\end{quote}

\subsection*{You'll want to know: \te{దుమ్ము}}
\textbf{\te{దుమ్ము}} — \emph{dummu} — \textquotedblleft{}dust\textquotedblright{}.

\begin{grammarlens}[title={What you've built}]
One more word for this chapter.
\end{grammarlens}

\subsection*{Your turn}
\begin{itemize}
\raggedright
  \item \emph{Say it:} \emph{dummu}
  \item \emph{Say it:} \emph{dummu}, once more
  \item \emph{Say it:} say \emph{dummu}
  \item \emph{Recall:} say the Telugu for the world, then the Telugu for an earthquake, then say \emph{dummu} again
\end{itemize}

\begin{tcolorbox}[breakable,colback=teal!4,colframe=teal!35!black,title={Before you move on}]
What does \te{దుమ్ము} mean? (\textquotedblleft{}Dust\textquotedblright{}.) Say it once more.
\end{tcolorbox}
\section[\texorpdfstring{\te{పొగ} (poga)}{పొగ (poga)}]{\te{పొగ} (poga) — smoke}

\label{lesson:TE-C175-poga}



\begin{quote}
Before the new one: say the Telugu for an earthquake, then the Telugu for dust.
\end{quote}

\subsection*{You'll want to know: \te{పొగ}}
\textbf{\te{పొగ}} — \emph{poga} — \textquotedblleft{}smoke\textquotedblright{}.

\begin{grammarlens}[title={What you've built}]
One more word for this chapter.
\end{grammarlens}

\subsection*{Your turn}
\begin{itemize}
\raggedright
  \item \emph{Say it:} \emph{poga}
  \item \emph{Say it:} \emph{poga}, once more
  \item \emph{Say it:} say \emph{poga}
  \item \emph{Recall:} say the Telugu for an earthquake, then the Telugu for dust, then say \emph{poga} again
\end{itemize}

\begin{tcolorbox}[breakable,colback=teal!4,colframe=teal!35!black,title={Before you move on}]
What does \te{పొగ} mean? (\textquotedblleft{}Smoke\textquotedblright{}.) Say it once more.
\end{tcolorbox}
